% --- Archon physics formalization source begin ---
% archon:physics
% archon:covers PhyXMiniProblems/problem_phyx_mini_0680.lean
% archon:source-report reports/phyx_mini/problem_phyx_mini_0680.source.json

\chapter{Physics problem phyx\_mini\_0680}
\label{ch:PhyXMiniProblems_problem_phyx_mini_0680}

\paragraph{Problem source.}
\#\# Physical scenario

The biggest stuffed animal in the world is a snake \textbackslash{}( 420 \textbackslash{}text\{ m\} \textbackslash{}) long, constructed by Norwegian children. Suppose the snake is laid out in a park as shown in figure, forming two straight sides of a \textbackslash{}( 105\textasciicircum{}\textbackslash{}circ \textbackslash{}) angle, with one side \textbackslash{}( 240 \textbackslash{}text\{ m\} \textbackslash{}) long. Olaf and Inge run a race they invent. Inge runs directly from the tail of the snake to its head, and Olaf starts from the same place at the same moment but runs along the snake.

\#\# Question

If Inge runs the race again at a constant speed of \textbackslash{}( 12.0 \textbackslash{}text\{ km/h\} \textbackslash{}), at what constant speed must Olaf run to reach the end of the snake at the same time as Inge

\#\# Answer choices

- A: A: \textbackslash{}( 16.0 \textbackslash{}text\{ km/h\} \textbackslash{})

- B: B: \textbackslash{}( 24.0 \textbackslash{}text\{ km/h\} \textbackslash{})

- C: C: \textbackslash{}( 15.0 \textbackslash{}text\{ km/h\} \textbackslash{})

- D: D: \textbackslash{}( 18.0 \textbackslash{}text\{ km/h\} \textbackslash{})

\#\# Figure caption (auxiliary; use the image as primary evidence)

The image depicts a stylized scene of a large, colorful snake in a park-like setting. The snake is composed of a series of patterned segments, each with a distinct design and color. The snake's body is extended in an L-shape across a grassy area.

Surrounding the snake, there are several trees depicted with circular, leafy tops. A winding path runs along the bottom left of the image, with a body of water visible next to it, creating a natural setting. The entire image has a playful, illustrative style.

\#\# Dataset metadata

Kinematics; Spatial Relation Reasoning, Multi-Formula Reasoning

\paragraph{Recorded answer/context.}
C: C: \textbackslash{}( 15.0 \textbackslash{}text\{ km/h\} \textbackslash{})

\paragraph{Figure/image path.}
/root/proposal\_for\_physic/science-mango/phyx\_mini\_run/phyx\_data/test\_image/680.png

\paragraph{Formalization target.}
create a compiling Lean file with sorry bodies at `PhyXMiniProblems/problem\_phyx\_mini\_0680.lean`. The Lean declarations must preserve the physical quantities, dimensions or dimensional roles, figure labels, governing-law hypotheses, and final relation expressed by this problem.
Use Mathlib/Physlib names found through LeanExplore where available. If a physics API is missing, introduce faithful local abstractions rather than scalar placeholder aliases.

\begin{theorem}[Physics formalization target]
\label{thm:physics:phyx_mini_0680:target}
\lean{PhyXMiniProblems.ProblemPhyXMini0680.problem_phyx_mini_0680}
\uses{def:physics:phyx-mini-0680:phyxminiproblems-problemphyxmini0680-speedreadout, def:physics:phyx-mini-0680:phyxminiproblems-problemphyxmini0680-radiansfromdegrees, def:physics:phyx-mini-0680:phyxminiproblems-problemphyxmini0680-snakeracesetup, def:physics:phyx-mini-0680:phyxminiproblems-problemphyxmini0680-matchessnakeracescenario, def:physics:phyx-mini-0680:phyxminiproblems-problemphyxmini0680-matchessuppliedsnakefigure, def:physics:phyx-mini-0680:phyxminiproblems-problemphyxmini0680-matchesproblemreadouts, def:physics:phyx-mini-0680:phyxminiproblems-problemphyxmini0680-satisfiestwostraightsidegeometry, def:physics:phyx-mini-0680:phyxminiproblems-problemphyxmini0680-hasphysicalraceparameters, def:physics:phyx-mini-0680:phyxminiproblems-problemphyxmini0680-satisfiesconstantspeedracekinematics, def:physics:phyx-mini-0680:phyxminiproblems-problemphyxmini0680-reachesheadatsametime, def:physics:phyx-mini-0680:phyxminiproblems-problemphyxmini0680-answerchoice, def:physics:phyx-mini-0680:phyxminiproblems-problemphyxmini0680-matchesdisplayedspeed}
The assigned autoformalize agent should translate this physics problem into Lean declarations in the covered file, with theorem and lemma proof bodies written as `by sorry`.
\end{theorem}
\begin{proof}
This is an autoformalization task, not a proof task. Produce statements that can later be proved by the physics prover without weakening the physical model.
\end{proof}

% --- Archon declaration topology begin ---
\section{Declaration topology}
\begin{definition}[Length Quantity]
  \label{def:physics:phyx-mini-0680:phyxminiproblems-problemphyxmini0680-lengthquantity}
  \lean{PhyXMiniProblems.ProblemPhyXMini0680.LengthQuantity}
  A nonnegative physical length, independent of any chosen unit system
\end{definition}
\begin{proof}
  This is the declared model definition of Length Quantity.
\end{proof}
\begin{definition}[Duration Quantity]
  \label{def:physics:phyx-mini-0680:phyxminiproblems-problemphyxmini0680-durationquantity}
  \lean{PhyXMiniProblems.ProblemPhyXMini0680.DurationQuantity}
  A nonnegative physical duration, independent of any chosen unit system
\end{definition}
\begin{proof}
  This is the declared model definition of Duration Quantity.
\end{proof}
\begin{definition}[Speed Quantity]
  \label{def:physics:phyx-mini-0680:phyxminiproblems-problemphyxmini0680-speedquantity}
  \lean{PhyXMiniProblems.ProblemPhyXMini0680.SpeedQuantity}
  A nonnegative physical speed from Physlib's unit-independent speed API
\end{definition}
\begin{proof}
  This is the declared model definition of Speed Quantity.
\end{proof}
\begin{definition}[length Readout]
  \label{def:physics:phyx-mini-0680:phyxminiproblems-problemphyxmini0680-lengthreadout}
  \lean{PhyXMiniProblems.ProblemPhyXMini0680.lengthReadout}
  \uses{def:physics:phyx-mini-0680:phyxminiproblems-problemphyxmini0680-lengthquantity}
  Read a physical length in a selected length unit
\end{definition}
\begin{proof}
  This is the declared model definition of length Readout.
\end{proof}
\begin{definition}[duration Readout]
  \label{def:physics:phyx-mini-0680:phyxminiproblems-problemphyxmini0680-durationreadout}
  \lean{PhyXMiniProblems.ProblemPhyXMini0680.durationReadout}
  \uses{def:physics:phyx-mini-0680:phyxminiproblems-problemphyxmini0680-durationquantity}
  Read a physical duration in a selected time unit
\end{definition}
\begin{proof}
  This is the declared model definition of duration Readout.
\end{proof}
\begin{definition}[speed Readout]
  \label{def:physics:phyx-mini-0680:phyxminiproblems-problemphyxmini0680-speedreadout}
  \lean{PhyXMiniProblems.ProblemPhyXMini0680.speedReadout}
  \uses{def:physics:phyx-mini-0680:phyxminiproblems-problemphyxmini0680-speedquantity}
  Read a physical speed in coherent selected length and time units
\end{definition}
\begin{proof}
  This is the declared model definition of speed Readout.
\end{proof}
\begin{definition}[radians From Degrees]
  \label{def:physics:phyx-mini-0680:phyxminiproblems-problemphyxmini0680-radiansfromdegrees}
  \lean{PhyXMiniProblems.ProblemPhyXMini0680.radiansFromDegrees}
  Convert the degree value printed in the problem to a radian readout
\end{definition}
\begin{proof}
  This is the declared model definition of radians From Degrees.
\end{proof}
\begin{definition}[Snake Vertex]
  \label{def:physics:phyx-mini-0680:phyxminiproblems-problemphyxmini0680-snakevertex}
  \lean{PhyXMiniProblems.ProblemPhyXMini0680.SnakeVertex}
  The three physically distinguished points along the laid-out snake
\end{definition}
\begin{proof}
  This is the declared model definition of Snake Vertex.
\end{proof}
\begin{definition}[Snake Side]
  \label{def:physics:phyx-mini-0680:phyxminiproblems-problemphyxmini0680-snakeside}
  \lean{PhyXMiniProblems.ProblemPhyXMini0680.SnakeSide}
  The two straight pieces of the snake, in tail-to-head order
\end{definition}
\begin{proof}
  This is the declared model definition of Snake Side.
\end{proof}
\begin{definition}[start Vertex]
  \label{def:physics:phyx-mini-0680:phyxminiproblems-problemphyxmini0680-snakeside-startvertex}
  \lean{PhyXMiniProblems.ProblemPhyXMini0680.SnakeSide.startVertex}
  \uses{def:physics:phyx-mini-0680:phyxminiproblems-problemphyxmini0680-snakevertex, def:physics:phyx-mini-0680:phyxminiproblems-problemphyxmini0680-snakeside}
  Initial endpoint of each oriented straight side
\end{definition}
\begin{proof}
  This is the declared model definition of start Vertex.
\end{proof}
\begin{definition}[finish Vertex]
  \label{def:physics:phyx-mini-0680:phyxminiproblems-problemphyxmini0680-snakeside-finishvertex}
  \lean{PhyXMiniProblems.ProblemPhyXMini0680.SnakeSide.finishVertex}
  \uses{def:physics:phyx-mini-0680:phyxminiproblems-problemphyxmini0680-snakevertex, def:physics:phyx-mini-0680:phyxminiproblems-problemphyxmini0680-snakeside}
  Final endpoint of each oriented straight side
\end{definition}
\begin{proof}
  This is the declared model definition of finish Vertex.
\end{proof}
\begin{definition}[Runner]
  \label{def:physics:phyx-mini-0680:phyxminiproblems-problemphyxmini0680-runner}
  \lean{PhyXMiniProblems.ProblemPhyXMini0680.Runner}
  The two named children who run the race
\end{definition}
\begin{proof}
  This is the declared model definition of Runner.
\end{proof}
\begin{definition}[Race Route]
  \label{def:physics:phyx-mini-0680:phyxminiproblems-problemphyxmini0680-raceroute}
  \lean{PhyXMiniProblems.ProblemPhyXMini0680.RaceRoute}
  The direct chord and the route following the snake itself
\end{definition}
\begin{proof}
  This is the declared model definition of Race Route.
\end{proof}
\begin{definition}[Park Figure Feature]
  \label{def:physics:phyx-mini-0680:phyxminiproblems-problemphyxmini0680-parkfigurefeature}
  \lean{PhyXMiniProblems.ProblemPhyXMini0680.ParkFigureFeature}
  Qualitative park features visibly present in the supplied raster
\end{definition}
\begin{proof}
  This is the declared model definition of Park Figure Feature.
\end{proof}
\begin{definition}[Snake Park Figure]
  \label{def:physics:phyx-mini-0680:phyxminiproblems-problemphyxmini0680-snakeparkfigure}
  \lean{PhyXMiniProblems.ProblemPhyXMini0680.SnakeParkFigure}
  \uses{def:physics:phyx-mini-0680:phyxminiproblems-problemphyxmini0680-snakevertex, def:physics:phyx-mini-0680:phyxminiproblems-problemphyxmini0680-parkfigurefeature}
  Planar physical layout associated with the primary figure. Coordinates are metre readouts; physical distances themselves remain LengthQuantity values
\end{definition}
\begin{proof}
  This is the declared model definition of Snake Park Figure.
\end{proof}
\begin{definition}[Snake Race Setup]
  \label{def:physics:phyx-mini-0680:phyxminiproblems-problemphyxmini0680-snakeracesetup}
  \lean{PhyXMiniProblems.ProblemPhyXMini0680.SnakeRaceSetup}
  \uses{def:physics:phyx-mini-0680:phyxminiproblems-problemphyxmini0680-lengthquantity, def:physics:phyx-mini-0680:phyxminiproblems-problemphyxmini0680-durationquantity, def:physics:phyx-mini-0680:phyxminiproblems-problemphyxmini0680-speedquantity, def:physics:phyx-mini-0680:phyxminiproblems-problemphyxmini0680-snakevertex, def:physics:phyx-mini-0680:phyxminiproblems-problemphyxmini0680-snakeside, def:physics:phyx-mini-0680:phyxminiproblems-problemphyxmini0680-runner, def:physics:phyx-mini-0680:phyxminiproblems-problemphyxmini0680-raceroute, def:physics:phyx-mini-0680:phyxminiproblems-problemphyxmini0680-snakeparkfigure}
  Independent observables and assignments in the race. In particular, Olaf's speed and the tail-to-head chord length are not defined from the answer choice
\end{definition}
\begin{proof}
  This is the declared model definition of Snake Race Setup.
\end{proof}
\begin{definition}[route Distance]
  \label{def:physics:phyx-mini-0680:phyxminiproblems-problemphyxmini0680-routedistance}
  \lean{PhyXMiniProblems.ProblemPhyXMini0680.routeDistance}
  \uses{def:physics:phyx-mini-0680:phyxminiproblems-problemphyxmini0680-lengthquantity, def:physics:phyx-mini-0680:phyxminiproblems-problemphyxmini0680-raceroute, def:physics:phyx-mini-0680:phyxminiproblems-problemphyxmini0680-snakeracesetup}
  The physical distance assigned to either possible race route
\end{definition}
\begin{proof}
  This is the declared model definition of route Distance.
\end{proof}
\begin{definition}[Matches Snake Race Scenario]
  \label{def:physics:phyx-mini-0680:phyxminiproblems-problemphyxmini0680-matchessnakeracescenario}
  \lean{PhyXMiniProblems.ProblemPhyXMini0680.MatchesSnakeRaceScenario}
  \uses{def:physics:phyx-mini-0680:phyxminiproblems-problemphyxmini0680-snakeracesetup}
  Runner identities, endpoints, routes, and constant-speed roles from the prose
\end{definition}
\begin{proof}
  This is the declared model definition of Matches Snake Race Scenario.
\end{proof}
\begin{definition}[Matches Supplied Snake Figure]
  \label{def:physics:phyx-mini-0680:phyxminiproblems-problemphyxmini0680-matchessuppliedsnakefigure}
  \lean{PhyXMiniProblems.ProblemPhyXMini0680.MatchesSuppliedSnakeFigure}
  \uses{def:physics:phyx-mini-0680:phyxminiproblems-problemphyxmini0680-radiansfromdegrees, def:physics:phyx-mini-0680:phyxminiproblems-problemphyxmini0680-snakeracesetup}
  Qualitative evidence and bend geometry read from 680.png. The angle is a given geometric datum, not a conclusion about either runner's speed
\end{definition}
\begin{proof}
  This is the declared model definition of Matches Supplied Snake Figure.
\end{proof}
\begin{definition}[Matches Problem Readouts]
  \label{def:physics:phyx-mini-0680:phyxminiproblems-problemphyxmini0680-matchesproblemreadouts}
  \lean{PhyXMiniProblems.ProblemPhyXMini0680.MatchesProblemReadouts}
  \uses{def:physics:phyx-mini-0680:phyxminiproblems-problemphyxmini0680-lengthreadout, def:physics:phyx-mini-0680:phyxminiproblems-problemphyxmini0680-speedreadout, def:physics:phyx-mini-0680:phyxminiproblems-problemphyxmini0680-snakeside, def:physics:phyx-mini-0680:phyxminiproblems-problemphyxmini0680-snakeracesetup}
  Numerical data stated in the problem. Which of the two sides has length 240 m is deliberately left open because the prose does not name that side. No value for Olaf's speed or the derived second side occurs here
\end{definition}
\begin{proof}
  This is the declared model definition of Matches Problem Readouts.
\end{proof}
\begin{definition}[Satisfies Two Straight Side Geometry]
  \label{def:physics:phyx-mini-0680:phyxminiproblems-problemphyxmini0680-satisfiestwostraightsidegeometry}
  \lean{PhyXMiniProblems.ProblemPhyXMini0680.SatisfiesTwoStraightSideGeometry}
  \uses{def:physics:phyx-mini-0680:phyxminiproblems-problemphyxmini0680-lengthreadout, def:physics:phyx-mini-0680:phyxminiproblems-problemphyxmini0680-snakeside, def:physics:phyx-mini-0680:phyxminiproblems-problemphyxmini0680-snakeracesetup}
  The two dimensionful side lengths realize the two straight planar segments; the direct route realizes the endpoint chord; and the route along the snake is the sum of its two sides. The cosine rule itself is supplied by Mathlib as EuclideanGeometry.law\_cos and is therefore not postulated as a custom law
\end{definition}
\begin{proof}
  This is the declared model definition of Satisfies Two Straight Side Geometry.
\end{proof}
\begin{definition}[Has Physical Race Parameters]
  \label{def:physics:phyx-mini-0680:phyxminiproblems-problemphyxmini0680-hasphysicalraceparameters}
  \lean{PhyXMiniProblems.ProblemPhyXMini0680.HasPhysicalRaceParameters}
  \uses{def:physics:phyx-mini-0680:phyxminiproblems-problemphyxmini0680-lengthreadout, def:physics:phyx-mini-0680:phyxminiproblems-problemphyxmini0680-durationreadout, def:physics:phyx-mini-0680:phyxminiproblems-problemphyxmini0680-speedreadout, def:physics:phyx-mini-0680:phyxminiproblems-problemphyxmini0680-snakeracesetup}
  Positivity conditions selecting the physical branch of lengths and times
\end{definition}
\begin{proof}
  This is the declared model definition of Has Physical Race Parameters.
\end{proof}
\begin{definition}[Satisfies Constant Speed Race Kinematics]
  \label{def:physics:phyx-mini-0680:phyxminiproblems-problemphyxmini0680-satisfiesconstantspeedracekinematics}
  \lean{PhyXMiniProblems.ProblemPhyXMini0680.SatisfiesConstantSpeedRaceKinematics}
  \uses{def:physics:phyx-mini-0680:phyxminiproblems-problemphyxmini0680-lengthreadout, def:physics:phyx-mini-0680:phyxminiproblems-problemphyxmini0680-durationreadout, def:physics:phyx-mini-0680:phyxminiproblems-problemphyxmini0680-speedreadout, def:physics:phyx-mini-0680:phyxminiproblems-problemphyxmini0680-snakeracesetup, def:physics:phyx-mini-0680:phyxminiproblems-problemphyxmini0680-routedistance}
  For each constant-speed runner, distance equals speed times elapsed time in every coherent choice of length and time units. This is the governing uniform-motion law and contains no numerical value for Olaf's speed
\end{definition}
\begin{proof}
  This is the declared model definition of Satisfies Constant Speed Race Kinematics.
\end{proof}
\begin{definition}[Reaches Head At Same Time]
  \label{def:physics:phyx-mini-0680:phyxminiproblems-problemphyxmini0680-reachesheadatsametime}
  \lean{PhyXMiniProblems.ProblemPhyXMini0680.ReachesHeadAtSameTime}
  \uses{def:physics:phyx-mini-0680:phyxminiproblems-problemphyxmini0680-durationreadout, def:physics:phyx-mini-0680:phyxminiproblems-problemphyxmini0680-snakeracesetup}
  The condition imposed by the question: after the simultaneous start, the two runners reach the common endpoint at the same time. It equates durations but does not constrain Olaf's speed to any proposed answer value
\end{definition}
\begin{proof}
  This is the declared model definition of Reaches Head At Same Time.
\end{proof}
\begin{definition}[Answer Choice]
  \label{def:physics:phyx-mini-0680:phyxminiproblems-problemphyxmini0680-answerchoice}
  \lean{PhyXMiniProblems.ProblemPhyXMini0680.AnswerChoice}
  Labels of the four speeds printed with the problem
\end{definition}
\begin{proof}
  This is the declared model definition of Answer Choice.
\end{proof}
\begin{definition}[displayed Speed Kilometers Per Hour]
  \label{def:physics:phyx-mini-0680:phyxminiproblems-problemphyxmini0680-displayedspeedkilometersperhour}
  \lean{PhyXMiniProblems.ProblemPhyXMini0680.displayedSpeedKilometersPerHour}
  \uses{def:physics:phyx-mini-0680:phyxminiproblems-problemphyxmini0680-answerchoice}
  Speed readout in kilometres per hour printed beside each answer label
\end{definition}
\begin{proof}
  This is the declared model definition of displayed Speed Kilometers Per Hour.
\end{proof}
\begin{definition}[Matches Displayed Speed]
  \label{def:physics:phyx-mini-0680:phyxminiproblems-problemphyxmini0680-matchesdisplayedspeed}
  \lean{PhyXMiniProblems.ProblemPhyXMini0680.MatchesDisplayedSpeed}
  \uses{def:physics:phyx-mini-0680:phyxminiproblems-problemphyxmini0680-speedquantity, def:physics:phyx-mini-0680:phyxminiproblems-problemphyxmini0680-speedreadout, def:physics:phyx-mini-0680:phyxminiproblems-problemphyxmini0680-answerchoice, def:physics:phyx-mini-0680:phyxminiproblems-problemphyxmini0680-displayedspeedkilometersperhour}
  Agreement with a speed displayed to the nearest 0.1 km/h
\end{definition}
\begin{proof}
  This is the declared model definition of Matches Displayed Speed.
\end{proof}
\begin{lemma}[direct Tail To Head Distance Meters]
  \label{lem:physics:phyx-mini-0680:phyxminiproblems-problemphyxmini0680-directtailtoheaddistancemeters}
  \lean{PhyXMiniProblems.ProblemPhyXMini0680.directTailToHeadDistanceMeters}
  \uses{def:physics:phyx-mini-0680:phyxminiproblems-problemphyxmini0680-lengthreadout, def:physics:phyx-mini-0680:phyxminiproblems-problemphyxmini0680-radiansfromdegrees, def:physics:phyx-mini-0680:phyxminiproblems-problemphyxmini0680-snakeracesetup, def:physics:phyx-mini-0680:phyxminiproblems-problemphyxmini0680-matchessuppliedsnakefigure, def:physics:phyx-mini-0680:phyxminiproblems-problemphyxmini0680-matchesproblemreadouts, def:physics:phyx-mini-0680:phyxminiproblems-problemphyxmini0680-satisfiestwostraightsidegeometry, def:physics:phyx-mini-0680:phyxminiproblems-problemphyxmini0680-hasphysicalraceparameters}
  The unstated side is 180 m; applying Mathlib's cosine rule to the three planar vertices gives the exact endpoint-chord readout. This derived geometry is a lemma, not a premise of the race theorem
\end{lemma}
\begin{proof}
  The conclusion follows from the governing relations and figure readouts named in the statement of direct Tail To Head Distance Meters.
\end{proof}
% --- Archon declaration topology end ---
% --- Archon physics formalization source end ---
